\documentclass[11pt,a4paper]{article}

\usepackage[T1]{fontenc}
\usepackage{lmodern}
\usepackage[margin=2.2cm]{geometry}
\usepackage{amsmath}
\usepackage{booktabs}
\usepackage{tabularx}
\usepackage{graphicx}
\usepackage{microtype}
\usepackage[hidelinks]{hyperref}

% Compact headings; each graded section title carries its criterion name and weight.
\makeatletter
\renewcommand\section{\@startsection{section}{1}{\z@}%
  {-2.0ex plus -0.5ex minus -0.2ex}{0.8ex plus 0.2ex}{\normalfont\large\bfseries}}
\renewcommand\subsection{\@startsection{subsection}{2}{\z@}%
  {-1.4ex plus -0.4ex minus -0.2ex}{0.4ex plus 0.1ex}{\normalfont\normalsize\bfseries}}
% Non-floating table with a numbered caption, used in the appendix so each table stays under its heading.
\newenvironment{fixedtable}%
  {\par\medskip\noindent\begin{minipage}{\linewidth}\def\@captype{table}\centering}%
  {\end{minipage}\par\medskip}
\makeatother

\setlength{\parindent}{0pt}
\setlength{\parskip}{0.5ex plus 0.2ex}
\setlength{\textfloatsep}{8pt plus 2pt minus 2pt}
\setlength{\floatsep}{8pt plus 2pt minus 2pt}
\setlength{\intextsep}{8pt plus 2pt minus 2pt}
\setlength{\abovecaptionskip}{3pt}
\setlength{\belowcaptionskip}{2pt}

\newcommand{\code}[1]{\texttt{#1}}
\newcommand{\Dsq}{D$^2$}

\begin{document}

\begin{center}
{\Large\bfseries Pricing 2025 claim cost in a developed and an emerging market}\\[4pt]
Team AIIAA (Yichuan Lin, Oussama El Hadi)\\
IAA AI-Assisted Hackathon 2026, 30 October 2026
\end{center}

\noindent We priced motor, home buildings and home contents for a developed market (Part 1) and an emerging market
(Part 2). The ranked Part 1 entry, S09, scored 0.13689 on the private board (rank 12 of 51); the team's final Part 1
method, a blend with a teammate's pipeline, scored 0.13718 as a late submission, which is not ranked. The Part 2
candidate P2-S15 scores 0.15807 on the public board (rank 18 of 38 on 30 October), against 0.14370 for our first
emerging-market file.

%======================================================================
\section{Pricing performance and adaptation (40\%)}

\subsection{Task, data and target}

Each 2025 policy row needs the ultimate cost of claims with a loss date in the row's period, net of excess and
recoveries and excluding declined claims, as a row total. Scoring is Tweedie \Dsq{} with power 1.5. Part 1 has
2,409,311 training rows for 2021 to 2024 and 163,150 claims valued at 31 December 2024; Part 2 has two training years,
2023 and 2024 (Table~\ref{tab:markets}).

\begin{table}[htbp]
\centering\small\setlength{\tabcolsep}{3.5pt}
\begin{tabular}{lrr@{\hspace{1em}}lrr}
\toprule
 & Developed & Emerging & & Developed & Emerging\\
\midrule
Test rows & 626,639 & 209,063 & New business in 2025 test rows & 15\% & 76\%\\
Motor exposure-years & 399,468 & 38,969 & Test policies seen in training & 85\% & 24\%\\
Motor claim frequency & 13.48\% & 18.09\% & Broker share, motor exposure & 25\% & 70\%\\
Motor mean severity (CU) & 3,862 & 2,250 & Missing values, where affected & 2.5\% & 12\%\\
\bottomrule
\end{tabular}
\caption{Developed against emerging market. Motor figures are for 2023-24, severity is reported cost, and the
missing-value shares are approximate.}
\label{tab:markets}
\end{table}

A claim attaches to the row of its policy whose period contains the loss date, at incurred minus recoveries.
Repudiated claims are excluded, and the 410 Part 1 claims that match no row are listed and left out; totals reconcile
with a difference of 0 (C005). The claims file is one snapshot, so instead of a triangle we fit a cost-weighted
reporting-lag curve per line on accident year 2021 and divide each row's reported cost by its completeness (C006).
Motor 2024 pure premium rises from 454 reported to 585 developed, inside the 564 to 606 range of 2021 to 2023. Bodily
injury is 43\% of motor cost and only 14\% of it is reported at 30 days. Claims are capped at the line's 99.5th
percentile and a per-line load adds the excess back (C014). Ten data issues, such as exposure stored in days and sums
insured inflated by 1,000, are handled in code at load time, and the raw files are never edited (C007,
Appendix~\ref{app:dq}).

\subsection{Part 1 model}

All models fit capped cost per unit of completeness-weighted exposure with a Tweedie loss of power 1.5 (C010): a
banded GLM, a spline GAM, LightGBM, XGBoost and CatBoost. Each is rescaled to its training total, because log-link
LightGBM ended at 0.83 of actual cost on its own training rows (C013), and a feature policy drops columns the trees
overfitted, such as motor \code{make\_model} (C026). Per-line blend weights from out-of-fold predictions are used only
if they beat equal weights on both folds (C021, C030). Half the fitted trend takes predictions to 2025 (motor 1.018,
buildings 1.095, contents 1.064; C019), and the final refit averages three seeds (C029). That is S09. Our teammate's
independent pipeline (OBLEMOS: a Poisson by Gamma GLM with BIC selection, a P-spline GAM, LightGBM and CatBoost,
averaged as Blend4) correlates 0.965 with S09. After the close we put Blend4 on S09's level per line and blended
$0.7\times\text{S09} + 0.3\times\text{Blend4}$; this late submission is the team's final Part 1 method.

\subsection{Adapting to the emerging market}

In the emerging market 76\% of the 2025 rows are new business, against 8\% of its training rows (C033), and the
construction codes, occupation text and age bands differ. The same code therefore runs with a market setting and a
shared schema (C034, C035). The adaptation took six steps, and each change was kept only if it passed the rule in
force at the time (Section~\ref{sec:trust}).

\begin{enumerate}\setlength{\itemsep}{1pt}\setlength{\parskip}{0pt}\setlength{\topsep}{2pt}
\item Emerging data alone (P2-S01). With about 500 claims per fit year on each property line, the GLM and GAM overfit
(buildings time-fold \code{d2\_cal} of $-10.5$ and $-30.7$).
\item Pooled training (C038, C040). Developed rows join every fit at weight $w$ times exposure, with \code{market} as a
feature, and the level is re-balanced on the emerging fit rows. Time-fold \code{d2\_cal} rose from 0.0762 to 0.0915 on
motor, 0.0410 to 0.1065 on buildings and 0.0779 to 0.1072 on contents; validation set $w = 1$ for buildings and $0.3$
for the other lines.
\item Level (C043). Buildings and contents keep the developed half trend (1.0948 and 1.0636), because the 2023 fire
year (buildings 2024 over 2023 is 0.63) rules out their own. Motor pure premium rose 10.7\% in 2024 against 1.75\%
applied, and six public-board probes (P2-S06 to P2-S11) moved motor to its own full trend, 1.1606. This is our one
decision taken mainly from the public board.
\item New business (C044, C047). The Part 2 audit found validation rows to be 92\% renewals, against 77\% new business
in the test (P2-F1), so views re-weighted to the test mix joined the gate. Contents became $0.35\times$ the pooled-0.3
GLM $+\,0.65\times$ the pooled-1 LightGBM (P2-S12). The 68,823 motor new-business rows take an equal 1/12 blend of GLM,
LightGBM, XGBoost and CatBoost from three training sets (pooled at weights 0.3 and 1, and new-business rows only),
because searched weights lost the guard (P2-S13). Motor renewal rows keep the four pooled-0.3 models at 0.25 each.
\item Teammate's GAM (C054). Refitted on the pooled data, the OBLEMOS Part 1 GAM beat the pooled CatBoost on the
buildings time fold (0.1257 against 0.1176). Buildings became $0.4\times$ that GAM $+\,0.6\times$ the pooled-1
CatBoost, at the P2-S13 buildings total (P2-S14).
\item Seed averaging (C057). P2-S15 averages the LightGBM, XGBoost and CatBoost test predictions over three seeds, as
S09 did, and keeps the P2-S13 buildings total (C058). Out-of-fold predictions do not change, so no fold gate applies;
the rule, fixed before the score was seen, was to keep it unless it fell more than 0.0005 below P2-S14.
\end{enumerate}

Of the 0.0144 public gain from P2-S01 to P2-S15, the developed half trend gave 0.0009, pooling 0.0073, the buildings
weight and motor level 0.0022 and the four steps of Table~\ref{tab:views} 0.0039 (Table~\ref{tab:results}).

\begin{table}[htbp]
\centering\small\setlength{\tabcolsep}{5pt}
\begin{tabular}{lrrrrrr}
\toprule
 & \multicolumn{2}{c}{Raw} & \multicolumn{2}{c}{Test mix} & \multicolumn{2}{c}{New business only}\\
\cmidrule(lr){2-3}\cmidrule(lr){4-5}\cmidrule(lr){6-7}
Adopted step & time & grouped & time & grouped & time & grouped\\
\midrule
P2-S12, contents blend & +0.0005 & +0.0032 & +0.0020 & +0.0041 & n/r & +0.0039\\
P2-S13, motor new-business rows & +0.0001 & +0.0009 & +0.0008 & +0.0050 & +0.0012 & +0.0032\\
P2-S14, buildings with the GAM & +0.0068 & +0.0018 & +0.0052 & +0.0058 & +0.0048 & +0.0068\\
P2-S15, three-seed refit & \multicolumn{6}{c}{out-of-fold predictions unchanged}\\
\bottomrule
\end{tabular}
\caption{Gain in \code{d2\_cal} on the six Part 2 validation views over the recipe each step replaced, on that line:
P2-S05 contents, P2-S12 motor and P2-S13 buildings (C044, C047, C054). The P2-S12 grouped gains were measured on the
folds its weights were searched on; its time-fold gains are out of search. n/r: not reported.}
\label{tab:views}
\end{table}

\begin{table}[htbp]
\centering\small
\begin{tabularx}{\textwidth}{llXrrl}
\toprule
 & File & Recipe & Public & Private & Rank\\
\midrule
Part 1 & S09 & Five models, guarded weights, half trend & 0.12863 & 0.13689 & 12 of 51\\
 & Blend4 & Teammate's pipeline alone & 0.12736 & 0.13526 & not ranked\\
 & S09 blend & $0.7\times$S09 $+\,0.3\times$Blend4, final Part 1 method & 0.12949 & 0.13718 & late, not ranked\\
\midrule
Part 2 & baseline & Constant average cost & 0.07319 & & \\
 & P2-S01 & Emerging data only & 0.14370 & & 8 of 10 (24 Oct)\\
 & P2-S04 & Pooled training, $w = 0.3$ on all lines & 0.15193 & & 9 of 12 (25 Oct)\\
 & P2-S10 & Buildings at $w = 1$, motor at its own trend & 0.15414 & & \\
 & P2-S13 & Contents blend, motor new-business segment & 0.15577 & & 23 of 35 (29 Oct)\\
 & P2-S14 & Teammate's GAM in buildings & 0.15781 & & \\
 & P2-S15 & P2-S14 with three-seed refit (final candidate) & 0.15807 & & 18 of 38 (30 Oct)\\
\bottomrule
\end{tabularx}
\caption{Kaggle Tweedie \Dsq{}. The Part 1 rank is the final private rank; Part 2 ranks are public-board team ranks on
the date shown, and Part 2 private scores follow the close on 31 October. On 30 October the public median was 0.15773
and the top score 0.17556.}
\label{tab:results}
\end{table}

Remaining uncertainty is mostly about level and new business. We assume unbiased case estimates, and that the accident
year 2021 reporting pattern holds in later years and in the emerging market; there the audit puts the overstatement of
the 2024 ultimate at 1.8\% for motor and 2.5\% for buildings (P2-F9). New-business predicted-to-actual ratios run 0.81
to 1.01 for motor, 1.30 to 2.23 for buildings and 0.74 to 0.84 for contents, on 37 to 697 claims per line and scheme,
too few to set a segment level. On a board of a third of the rows, one line's optimal level has a sampling standard
deviation of 4.7\% for motor (audit H1), and in Part 1 the private board reversed the public preference for half trend
over no trend. Appendix~\ref{app:lim} lists the other limitations.

\subsection{Office premium for the 20 reference policies}

$R$ is the P2-S15 prediction for the row, already a row total. From the emerging column of \code{expenses.parquet},
catastrophe reinsurance is $\text{Cat} = 0.12\,R$ on buildings and contents and 0 on motor (its unit is a share of
property risk premium, and we count contents as property). The fixed expense $F = 85$~CU is charged in full per policy.
The share-of-premium items add to $k = 0.15$ (variable expense) $+$ commission (0.20 broker, 0 direct) $+\ 0.08$
(premium tax) $+\ 0.08$ (target profit); none of the 20 policies uses the comparison-site channel, whose value is blank.
There is no minimum or maximum premium.
\[
P = \frac{R + \text{Cat} + F}{1-k},\qquad 1-k = 0.49 \text{ (broker)},\quad 0.69 \text{ (direct)}.
\]
For B00229821, $387.92 + 46.55 + 85.00 = 519.47$ and $519.47/0.49 = 1{,}060.14$. Table~\ref{tab:op} gives all 20
policies and Appendix~\ref{app:op} splits each $P$ by item.

\begin{table}[htbp]
\centering\small\setlength{\tabcolsep}{4pt}
\begin{tabular}{lrrrr@{\hspace{1.4em}}lrrrr}
\toprule
row\_id & $R$ & Cat & $P$ & WP & row\_id & $R$ & Cat & $P$ & WP\\
\midrule
M00118569$^{\mathrm{d,r}}$ & 922.76 & 0.00 & 1,460.52 & 1,095.48 & B00218282 & 391.04 & 46.92 & 1,067.27 & 230.80\\
M00129413 & 212.69 & 0.00 & 607.53 & 354.15 & B00223369 & 515.20 & 61.82 & 1,351.08 & 204.56\\
M00153851 & 368.16 & 0.00 & 924.81 & 414.72 & B00229821 & 387.92 & 46.55 & 1,060.14 & 285.60\\
M00161473 & 1,135.24 & 0.00 & 2,490.28 & 1,131.78 & B00250679 & 38.71 & 4.65 & 261.95 & 180.72\\
M00167496 & 777.40 & 0.00 & 1,760.00 & 463.95 & C00259568$^{\mathrm{r}}$ & 40.88 & 4.91 & 266.92 & 83.18\\
M00176777 & 1,072.30 & 0.00 & 2,361.85 & 443.61 & C00270331 & 37.45 & 4.49 & 259.07 & 140.29\\
M00185311$^{\mathrm{d}}$ & 1,105.96 & 0.00 & 1,726.03 & 872.31 & C00283470 & 86.59 & 10.39 & 371.39 & 117.97\\
B00215707 & 342.12 & 41.05 & 955.46 & 272.44 & C00286287 & 162.14 & 19.46 & 544.08 & 269.23\\
B00215809$^{\mathrm{d}}$ & 480.49 & 57.66 & 903.11 & 237.22 & C00302985 & 60.81 & 7.30 & 312.46 & 115.88\\
B00217584 & 825.49 & 99.06 & 2,060.29 & 323.36 & C00314425 & 5.97 & 0.72 & 187.11 & 102.87\\
\bottomrule
\end{tabular}
\caption{Office premium build-up in CU, with $F = 85.00$ on every row. Row ids starting M, B and C are motor,
buildings and contents. All rows are broker new business ($1-k = 0.49$) except d: direct ($1-k = 0.69$) and r:
renewal. WP is the written premium under the current tariff. Source: \code{scripts/part2\_office\_premium.py
-{}-predictions work/p2\_s15\_seeds3.csv} (C060).}
\label{tab:op}
\end{table}

The 20 rows sum to $R = 8{,}969.31$ and $P = 20{,}931.34$, against 7,340.12 of written premium, which is the current
tariff and not a target for our 2025 cost. $P$ exceeds written premium on every row, by 1.33 to 6.60 times (median 2.92), mostly through the loadings: a broker row with $R = 0$ would still cost
173.47, and $F$, Cat and the shares make up 36\% to 97\% of $P$. $R$ alone exceeds written premium on 10 rows, 6
buildings and 4 motor, all new business; our buildings new-business predictions ran 1.30 times actual on the grouped
folds (77 claims), so part of that gap may be our own overstatement. Written premium does not scale with exposure, so
$F$ dominates the three short rows. Pro-rating $F$ by exposure would lower $P$ from 607.53 to 491.56 for M00129413
(exposure 0.33), from 261.95 to 107.49 for B00250679 (0.11) and from 187.11 to 25.52 for C00314425 (0.07). These
conventions are sign-off item R15.

%======================================================================
\section{Effective use of AI (25\%)}

\subsection{How the AI worked}

We used Claude Code in three roles. A coordinator (claude-fable-5-1) planned each step, wrote the decision rule before
a run and proposed what to submit. Workers (claude-opus-5-5 and claude-sonnet-5-5) wrote a failing test, then the code,
ran the experiment and logged the change in \code{docs/CHANGELOG.md}. Read-only auditors trained nothing and backed each
finding with a script: 17 scripts behind B1 to B8 in Part 1 and 27 checks behind P2-F1 to P2-F10 in Part 2.
\code{docs/AI\_LOG.md} (A01 to A32) records for every step what the AI produced, how it was checked and whether it was
accepted, modified or rejected. Our teammate ran his Part 2 experiments in his own Claude session (C048 to C056, A27).

The AI worked in every stage. It profiled all 11 files against the sample submission (A02), counted each data issue
with a script (A03), scanned every column for leakage (A06) and wrote features and models test-first (A07). In the first full run it traced an XGBoost crash to a boolean column with missing values,
and a 9 to 18\% under-prediction of the trees to the loss of balance in log-link Tweedie boosting, testing each cause
before changing code (A08). Its vectorised weight search covers 888,030 weight vectors for eight candidates, and the
threaded version cut the motor search from 1,277 to 306 seconds with the same weights (C041).

\subsection{Where a check corrected or stopped the AI}

AI output reached a submission only after a check. Table~\ref{tab:ai} lists proposals that failed one; the checks came
from other agents, from tests and from the data.

\begin{table}[htbp]
\centering\small
\begin{tabularx}{\textwidth}{>{\raggedright\arraybackslash}p{4.6cm}>{\raggedright\arraybackslash}X>{\raggedright\arraybackslash}p{3.9cm}}
\toprule
AI proposal & Check & Outcome\\
\midrule
Per-peril reporting curves (coordinator, C022) & Worker compared them with the pooled curve: largest difference 5e-16 & No effect; kept as an input (A09)\\
GLM market interactions in a 9-model blend, +0.0024, above the rule (C042) & Worse alone on all six contents folds; new-business level up 21\% with no test of it & Rejected by the coordinator (A23)\\
Buildings renewal segment that passed the gate (C047) & Its gain came from a new-business-only GLM applied to renewal rows & Rejected (A26)\\
Motor control run, better on every view (C047) & It was run after the result was seen & Diagnostic only, not submitted (A26)\\
Ridge GLM in the transfer model (C048) & Shrinkage test: a penalty of 1e9 must return the base model & Scaling bug fixed before any result (A27)\\
Note that a teammate's Part 1 file scored 0.1356 in public & Kaggle listing with labelled columns: 0.13562 was its private score & Note corrected; S09 led on both boards (A28)\\
\bottomrule
\end{tabularx}
\caption{AI proposals that a check changed or stopped (\code{docs/AI\_LOG.md}, \code{docs/CHANGELOG.md}).}
\label{tab:ai}
\end{table}

\subsection{What the humans decided}

People made the judgement calls. The team lead approved items R01 to R11 of the sign-off sheet (cleaning rules,
development, cap, trend, feature drops, Part 2 schema, borrowed curves and pooled design) on 25 October in a single
reply without item-level reasons, and the sheet records that. R12 to R16 (the motor level, the test-mix gate, the
new-business segment, the office-premium conventions and P2-S15) were signed the same way on 30 October, so all 16 are
signed. The team asked for XGBoost and the GAM (A07, C011),
built the OBLEMOS pipeline, chose the $0.7/0.3$ blend as the Part 1 method and P2-S15 as the Part 2 final with P2-S14 as
the hedge. After the misread score in Table~\ref{tab:ai}, the Part 2 selection compared scores from one board, with
the six views, on files the repository rebuilds.

%======================================================================
\section{Trust, validation and guardrails (20\%)}
\label{sec:trust}

\subsection{Validation design}

No split is random and the year is not a feature. Part 1 uses two expanding-window folds that score 2023 and 2024.
Part 2 has a time fold (train 2023, score 2024) and a grouped 5-fold split on \code{policy\_id}, both scored on emerging
rows (C036). Our validation score \code{d2\_cal} rescales each fold to its actual total, so it measures ranking, and trees
early-stop on a random 10\% of the fitting rows, never on the scored fold (C027). Each Part 2 candidate is scored on six
views: the time fold and the grouped mean, each raw, re-weighted so that new business carries its share of 2025 test
exposure (about 0.77) and on new-business rows only (C044). The re-weighted view
tracks the board: the P2-S03 to P2-S04 step was +0.0154 raw, +0.0051 re-weighted and +0.0050 in public. Weights are
still searched on raw folds: 37 to 697 new-business claims can reject a change but not tune one.

\subsection{Adoption rules}

Rules were fixed before the runs they judged. A Part 1 change needed a mean gain of at least 0.002 with no fold lost
(C018). Searched blend weights must beat equal weights on every scheme, enforced in code (C030). A
Part 2 change must not lose on the raw time fold, the raw grouped mean or the test-mix grouped mean, and must gain
0.002 on the raw mean or the test-mix grouped mean (C044). After a contents candidate with negative time-fold test-mix
and new-business views lost 0.00017 in public, we also ruled out any loss on those two views (C053). Gaps under 0.0005
are ties, won by the simpler option. Among the Part 2 blend candidates tried in public, only those that gained on all
six views raised the score.

\subsection{Data, code and reproduction}

A SHA256 manifest, read-only file attributes and a test guard every raw byte (C007), and the developed files in the
Part 2 release matched the Part 1 manifest (C033). The suite has 341 pytest cases in 24 files, and modules are written
test-first. The blend step rebuilds S09 byte for byte from the stored model predictions (same SHA256 as the upload; a
full refit differs slightly between runs, as the tree refits are not bit-reproducible), and P2-S14, rebuilt by two
routes, differs from the upload by at most 5.1e-10 (C058). P2-S15 keeps the row ids, row order and buildings total of
P2-S14 (correlation 0.9991).

\subsection{Audits and human sign-off}

The Part 1 audit's findings on \code{make\_model}, \code{branch\_id} and early stopping became C026 and C027; its trend
double-counting finding (B4) was probed and not adopted. The Part 2 audit changed our validation (P2-F1) and caught a
diagnostic that read uncorrected exposure (P2-F8). \code{docs/HUMAN\_REVIEW.md} lists 16 decisions (R01 to R16) with
their evidence by file and line.

%======================================================================
\section{Communication and documentation (15\%)}

\code{README.md} gives the task, layout, pinned environment, commands, pipeline, governance and known limitations. \code{docs/CHANGELOG.md} has 61 entries (C000 to C060), each with its reason, evidence and code
location, and code comments cite the same ids. \code{docs/SUBMISSIONS.md} lists every upload with its recipe,
validation and scores. \code{docs/AI\_LOG.md}, \code{docs/HUMAN\_REVIEW.md} and the two audit files hold the AI record,
\code{docs/references/} archives the Kaggle pages, and \code{docs/part2\_validation/} keeps the Part 2 fold scores. The
evidence board in \code{docs/dashboard/} puts the main figures on one page, and the five-minute video follows
\code{docs/VIDEO\_SCRIPT.md}. The sign-off sheet and team notes are in Chinese; code and logs are in English.

A reviewer rebuilds each file as follows. \code{reproduce.ps1} runs from the raw data to S09 (the blend step from stored predictions matches the
upload's SHA256; a full refit differs slightly between runs) and with \code{-Part2} to P2-S13. \code{docs/PART2\_OBLEMOS.md} gives the two commands for P2-S14, which
\code{blend\_runs.py -{}-write-from} with \code{-{}-level-from} also builds (C058). P2-S15 is that write step after a
three-seed refit (\code{oof\_predictions.py -{}-refit-only -{}-seeds 3}), read with \code{-{}-pred-tag s3} (C057, C059). The
Part 1 blend comes from \code{oblemos\_part1/blend\_with\_S09.py}.

%======================================================================
\clearpage
\appendix
\begin{center}
{\Large\bfseries Appendix, not part of the five graded pages}
\end{center}

\section{Data-quality issues and treatment}
\label{app:dq}

\begin{fixedtable}
\centering\small
\begin{tabularx}{\textwidth}{>{\raggedright\arraybackslash}p{4.9cm}>{\raggedright\arraybackslash}p{4.2cm}>{\raggedright\arraybackslash}X}
\toprule
Issue & Size & Treatment\\
\midrule
Q1 \code{exposure} holds day counts or a 360-day basis (1.0139 = 365/360) & 5,723 motor, 4,176 buildings, 4,305 contents rows & Recomputed from period dates on a 365-day year (C001)\\
Q2 motor \code{customer\_id} lost leading zeros & 98,110 training and 25,554 test rows & Padded to 8 characters; 100\% join to customers (C002)\\
Q3 \code{sum\_insured} inflated by 1,000 & Values above 5m, including 370 buildings and 360 contents test rows & Divided by 1,000 and flagged; the same policies differ by exactly 1,000 across years (C003)\\
Q4 \code{construction\_type} case and codes T/W & 13,250 policies switch from T to W & Lower-cased; T and W mapped to timber (C004)\\
Q5 claims matching no policy row & 294 motor, 67 buildings, 49 contents & Excluded from the target and listed; reconciliation difference 0 (C005)\\
Q6 \code{customer\_service\_contacts} & About 3.0 on claimed rows against 0.41 on unclaimed rows & Dropped as leakage (C008)\\
Q7 about 2.5\% random missing values & Many columns & Not imputed for tree models; the GAM uses median imputation plus missing indicators (C011)\\
Q8 \code{telematics\_score} empty for 87\% of motor rows & Structural: no device & Not imputed\\
Q9 \code{renewal\_decision} only in training & All rows & Dropped; known only after the row ends\\
Q10 training and test drift & Median \code{renewal\_count} 6 to 3, \code{vehicle\_age} 7 to 9, \code{ncb\_years} 7 to 8 & Year is not a feature; the audit found the drift internally consistent (\code{renewal\_count} +1.0 per year)\\
\bottomrule
\end{tabularx}
\caption{Data-quality issues in the developed market and their treatment (\code{docs/DATA\_GUIDE.md}, \code{docs/CHANGELOG.md}).}
\end{fixedtable}

Part 2 re-check (C033): Q1 recurs in the emerging training tables only (253 motor rows); Q2 is absent; Q3 recurs and
the 5m threshold still separates; Q4 changes to the codes BRK, STN, WOD and MIX (C035); 36 claims match no row; missing
values rise to about 12\%; the Q8 columns are absent. \code{customer\_service\_contacts} again separates claimed rows (3.4
against 0.42 contacts) and stays out.

\section{All submissions}

\begin{fixedtable}
\centering\small
\begin{tabularx}{\textwidth}{lXrr}
\toprule
ID & Recipe & Public & Private\\
\midrule
S01 & Equal blend of GLM, GAM, LightGBM, XGBoost; full trend & 0.12453 & 0.13347\\
S02 & S01 with no trend & 0.12376 & 0.13466\\
S03 & S01 with half trend & 0.12464 & 0.13456\\
S04 & Tuned LightGBM and XGBoost, equal weights, half trend & 0.12477 & 0.13279\\
S05 & Equal blend of GLM, GAM and tuned trees & 0.12464 & 0.13477\\
S06 & Per-line weights from out-of-fold predictions (C021) & 0.12579 & 0.13505\\
S07 & Feature policy, inner early stopping, CatBoost, guarded weights & 0.12851 & 0.13693\\
S08 & S07 with a lower trend on buildings and contents (probe) & 0.12820 & 0.13703\\
S09 & S07 recipe with seed-averaged refit (ranked final, rank 12 of 51) & 0.12863 & 0.13689\\
Blend4 & Teammate's pipeline alone (not ranked) & 0.12736 & 0.13526\\
Late & $0.5\times$S09 $+\,0.5\times$Blend4 & 0.12971 & 0.13707\\
Late & $0.7\times$S09 $+\,0.3\times$Blend4 (final Part 1 method) & 0.12949 & 0.13718\\
\bottomrule
\end{tabularx}
\caption{Part 1 submissions. Private scores come from the Kaggle listing after the close. Late submissions are not
ranked; the scored $0.7/0.3$ file was the teammate's upload, and the repository holds a copy rebuilt by
\code{oblemos\_part1/blend\_with\_S09.py}.}
\end{fixedtable}

\begin{fixedtable}
\centering\small
\begin{tabularx}{\textwidth}{lXr}
\toprule
ID & Recipe & Public\\
\midrule
P2-S01 & Emerging data only, no trend & 0.14370\\
P2-S02 & P2-S01 with developed half trend & 0.14460\\
P2-S03 & Contents pooled & 0.14693\\
P2-S04 & All lines pooled, $w = 0.3$ & 0.15193\\
P2-S05 & P2-S04 with buildings at $w = 1.0$ & 0.15123\\
P2-S06 & Level probe, all lines $\times\,1.10$ & 0.15206\\
P2-S07 & Level probe, motor $\times\,1.10$ & 0.15362\\
P2-S08 & Level probe, buildings $\times\,1.10$ & 0.15029\\
P2-S09 & Motor at the square root of its own trend, 1.0773 & 0.15284\\
P2-S10 & Motor at its full own trend, 1.1606 & 0.15414\\
P2-S11 & P2-S10 uplift on motor new business only (tie, not adopted) & 0.15434\\
P2-S12 & P2-S10 with the contents blend across weight runs & 0.15474\\
P2-S13 & P2-S12 with the motor new-business segment blend & 0.15577\\
P2-S14 & P2-S13 with buildings $= 0.4\times$ teammate GAM $+\,0.6\times$ pooled CatBoost, at the P2-S13 buildings total & 0.15781\\
P2-S15 & P2-S14 with three-seed averaged LightGBM, XGBoost and CatBoost (final candidate; hedge P2-S14) & 0.15807\\
\bottomrule
\end{tabularx}
\caption{Part 2 submissions, public scores (\code{docs/SUBMISSIONS.md}). Five further probes uploaded by our teammate
on P2-S12 and P2-S13 scored 0.15483 to 0.15563 and none was adopted.}
\end{fixedtable}

\section{Part 2 validation detail}

\begin{fixedtable}
\centering\small
\begin{tabular}{lllll}
\toprule
Line & Emerging only & Pooled, $w = 0.3$ & Pooled, $w = 1.0$ & Adopted $w$\\
\midrule
Motor & 0.0762 / 0.0677 & 0.0915 / 0.0740 & 0.0908 / 0.0730 & 0.3\\
Buildings & 0.0410 / 0.0732 & 0.1065 / 0.1132 & 0.1169 / 0.1180 & 1.0\\
Contents & 0.0779 / 0.0934 & 0.1072 / 0.1100 & 0.1069 / 0.1128 & 0.3\\
\bottomrule
\end{tabular}
\caption{Emerging-only against pooled training, \code{d2\_cal} (time fold / grouped mean).}
\end{fixedtable}

The GAM left the pooled runs after it exceeded the memory cap on 756k rows. Tree models needed level re-balancing on
emerging rows, with factors of 1.08 to 1.26 at weight 0.3 on the property lines and 1.01 to 1.10 at weight 1.0.
P2-S05 (buildings at $w = 1.0$) scored 0.15123 in public against 0.15193 for P2-S04, so the board disagreed with
validation; we kept the validation rule (C040).

\begin{fixedtable}
\centering\small
\begin{tabular}{lrrrr}
\toprule
Recipe & Time & Grouped & Grouped test-mix & Public\\
\midrule
P2-S01 & 0.10971 & 0.12306 & 0.13006 & 0.14370\\
P2-S03 & 0.11546 & 0.12623 & 0.13421 & 0.14693\\
P2-S04 & 0.14179 & 0.14165 & 0.13934 & 0.15193\\
P2-S05 & 0.14486 & 0.14311 & 0.13973 & 0.15123\\
\bottomrule
\end{tabular}
\caption{Audit evidence for P2-F1: Kaggle-like \Dsq{} of the three lines stacked, each line sum-calibrated.}
\end{fixedtable}

\begin{fixedtable}
\centering\small
\begin{tabular}{llll}
\toprule
Line & Time fold, new / renewal & Grouped, new / renewal & New-business claims, time / grouped\\
\midrule
Motor & 1.012 / 0.933 & 0.812 / 1.025 & 382 / 697\\
Buildings & 2.229 / 1.607 & 1.296 / 0.981 & 37 / 77\\
Contents & 0.737 / 1.111 & 0.837 / 1.018 & 65 / 102\\
\bottomrule
\end{tabular}
\caption{Predicted-to-actual on reported cost for the P2-S05 blend, by segment. Source:
\code{docs/part2\_new\_business\_check.md}.}
\end{fixedtable}

\begin{fixedtable}
\centering\small
\begin{tabular}{lrrrr}
\toprule
Line & Days & Emerging 2023 & Developed 2023, same truncation & Developed 2021, used\\
\midrule
Motor & 30 & 0.603 & 0.579 & 0.556\\
Motor & 90 & 0.774 & 0.717 & 0.723\\
Motor & 365 & 0.961 & 0.970 & 0.922\\
Buildings & 30 & 0.913 & 0.829 & 0.857\\
Buildings & 90 & 0.952 & 0.927 & 0.933\\
Buildings & 365 & 0.998 & 0.992 & 0.992\\
Contents & 30 & 0.961 & 0.947 & 0.949\\
Contents & 90 & 0.991 & 0.995 & 0.989\\
Contents & 365 & 0.993 & 0.996 & 0.990\\
\bottomrule
\end{tabular}
\caption{Cost-weighted cumulative share reported (\code{docs/part2\_emerging\_target\_summary.txt}, C037).}
\end{fixedtable}

The peril mix differs most on buildings: fire is 0.454 of emerging cost against 0.237 in developed accident year 2021.
Mean completeness of the 2024 motor rows is 0.7672 with the developed peril shares and 0.8202 with the emerging ones.

\section{Office premium split by item}
\label{app:op}

\begin{fixedtable}
\centering\small
\begin{tabular}{lrrrrrr}
\toprule
row\_id & $R$ + Cat + $F$ & Variable expense & Commission & Premium tax & Profit & $P$\\
\midrule
M00118569 & 1,007.76 & 219.08 & 0.00 & 116.84 & 116.84 & 1,460.52\\
M00129413 & 297.69 & 91.13 & 121.51 & 48.60 & 48.60 & 607.53\\
M00153851 & 453.16 & 138.72 & 184.96 & 73.99 & 73.99 & 924.81\\
M00161473 & 1,220.24 & 373.54 & 498.06 & 199.22 & 199.22 & 2,490.28\\
M00167496 & 862.40 & 264.00 & 352.00 & 140.80 & 140.80 & 1,760.00\\
M00176777 & 1,157.30 & 354.28 & 472.37 & 188.95 & 188.95 & 2,361.85\\
M00185311 & 1,190.96 & 258.90 & 0.00 & 138.08 & 138.08 & 1,726.03\\
B00215707 & 468.17 & 143.32 & 191.09 & 76.44 & 76.44 & 955.46\\
B00215809 & 623.14 & 135.47 & 0.00 & 72.25 & 72.25 & 903.11\\
B00217584 & 1,009.54 & 309.04 & 412.06 & 164.82 & 164.82 & 2,060.29\\
B00218282 & 522.96 & 160.09 & 213.45 & 85.38 & 85.38 & 1,067.27\\
B00223369 & 662.03 & 202.66 & 270.22 & 108.09 & 108.09 & 1,351.08\\
B00229821 & 519.47 & 159.02 & 212.03 & 84.81 & 84.81 & 1,060.14\\
B00250679 & 128.36 & 39.29 & 52.39 & 20.96 & 20.96 & 261.95\\
C00259568 & 130.79 & 40.04 & 53.38 & 21.35 & 21.35 & 266.92\\
C00270331 & 126.94 & 38.86 & 51.81 & 20.73 & 20.73 & 259.07\\
C00283470 & 181.98 & 55.71 & 74.28 & 29.71 & 29.71 & 371.39\\
C00286287 & 266.60 & 81.61 & 108.82 & 43.53 & 43.53 & 544.08\\
C00302985 & 153.10 & 46.87 & 62.49 & 25.00 & 25.00 & 312.46\\
C00314425 & 91.68 & 28.07 & 37.42 & 14.97 & 14.97 & 187.11\\
\bottomrule
\end{tabular}
\caption{Each share-of-premium item in CU; every row adds up to $P$. Exact exposures and the full table come from
\code{scripts/part2\_office\_premium.py -{}-predictions work/p2\_s15\_seeds3.csv}.}
\end{fixedtable}

\section{Part 1 target and validation detail}

Repudiated claims (1,711, all with zero amounts) are excluded (C005), and 20.7\% of 2024 motor claims were still open
at 31 December 2024. The per-line caps are 51,585, 147,445 and 28,449 CU for motor, buildings and contents in Part 1
and 27,578, 220,808 and 28,312 CU in the emerging market (C014, C037); the capped share of cost is 8.8\%, 4.2\% and
2.1\% in Part 1 and 5.0\%, 4.9\% and 2.1\% in the emerging market. Motor perils other than bodily injury are 76\% to
99\% reported at 30 days.

Rescaling each model to its training total raised motor validation \Dsq{} from 0.066 to 0.073 (C013). Per-line tuning was
adopted only where the mean gain was at least 0.002: motor LightGBM gained 0.0075, motor XGBoost 0.0031 and contents
XGBoost 0.0024 (C015, C018). The feature policy (C026) raised mean LightGBM \code{d2\_cal} from 0.0832 to 0.0887 on
motor, from 0.1094 to 0.1134 on buildings and from 0.0891 to 0.0936 on contents. Motor searched weights lost to equal
weights on the 2024 fold (0.09031 against 0.09071), so motor uses equal weights (C021, C030). The S09 blends scored a
mean two-fold \code{d2\_cal} of 0.0919 for motor, 0.1144 for buildings and 0.0945 for contents.

Rejected Part 1 alternatives: per-peril frequency by severity scored 0.0801 against 0.0893 on motor, 0.1045 against
0.1094 on buildings and 0.0879 against 0.0911 on contents, with 2024 bodily-injury counts over-predicted by 1.43 (C023).
An out-of-fold target encoding of \code{make\_model} carried weak signal (code relativities correlated 0.18 across years,
against $-0.05$ for shuffled codes), and the equal blend gained 0.0003 on the mean but fell on 2024 (C031). Re-tuning
the trees on the final feature set gave motor LightGBM +0.0027 on the mean but $-0.0008$ on 2024, so nothing was adopted
(C032).

\section{Assumptions and limitations}
\label{app:lim}

Both parts assume that open-claim case estimates are unbiased and that the accident year 2021 reporting pattern holds
for later years and for the emerging market. Part 1 also assumes that the construction codes T and W both mean timber
(C004).

The emerging market's two years give one time fold, with about 500 claims per year on each property line, and 2024 is
also the least developed year. The reporting curves are borrowed, nothing beyond 24 months can be checked, and 20\% of
2024 emerging motor claims are open. Buildings 2023 was a fire year (49 fire claims for 3.22m CU against 21 for 1.13m
in 2024). The large-loss load is uniform within a line, and the emerging caps rest on 36, 6 and 6 claims above the cap.

Part of the level rests on public-board probes: the developed half trend (P2-S02, +0.0009, below our 0.002 rule) and
the motor level. The motor gain (+0.0029) is larger, and the model trained on 2023 had also under-predicted 2024
motor by 6\%. The risk is that a level read from about a third of the rows does not hold on the rest. In Part 1 the S08 probe of audit finding B4 flipped sign between the boards (public
$-0.0003$, private $+0.0001$).

Part 1 audit items left open: the trend may count the year level twice (B4; residual trend 1.017, 1.048 and 1.004
against 1.018, 1.095 and 1.064 applied). The 2025 motor large-loss load includes under-reported 2024 bodily injury and
may be 1.4\% low (B5). Recoveries appear only on closed claims, so open 2024 own-damage cost is about 1.4\% overstated
(B6). Prior-claim history is developed in cost but not in claim counts (B7). Open-claim reserving for bodily injury
cannot be verified from one snapshot, and the customer snapshot date is unknown, so early training rows may see later
information. Buildings and contents cost above the cap concentrates in the top risk decile while the load is uniform,
and the GAM uses a single penalty (C011). In the motor new-business segment, a control run made after the result was
seen scored higher on every view; choosing it would have been a post-hoc selection, so it stays a diagnostic.

\end{document}
